\documentclass[book,paper=a4,fontsize=10pt,jafontscale=0.92]{jlreq}

\usepackage{newtxtext,newtxmath}
\usepackage{graphicx}
\usepackage[usenames]{xcolor}
\usepackage{amsmath}
\usepackage{bm}
\usepackage{url}

\graphicspath{{./figure/}}

% スタイルファイルを1つ選択する(先頭の % を外す) %%%%%%%%%%%%%%%%%%%%%%%%%%%

% 知能・機械工学課程(学部)/知能・機械工学専攻(大学院)
%\usepackage{am-thesis} % 正式提出用
%\usepackage[check]{am-thesis} % チェック用

%人間システム工学科(学部)/人間システム工学専攻(大学院)
\usepackage{hsi-thesis} % 正式提出用
%\usepackage[check]{hsi-thesis} % チェック用

% 以下を適切に設定する %%%%%%%%%%%%%%%%%%%%%%%%%%%%%%%%%%%%%%%%%%%%%%%%%%%%%

% 修士論文/卒業論文の選択

\thesistype{M} % 修士論文はM 卒業論文はB

\title{VRコンテンツ体験時における生体情報に\\基づくリアルタイム感情推定フレームワーク} % あなたの論文のタイトルをここに書く
\author{趙 聖化} % 名字と名前の間は *半角* スペースを入れる．
\studentnumber{27020856} % 8桁で書く．

\syear{2026} % 提出日
\smonth{3}
\sday{1} % 日は出力されない

\mentor{井村 誠孝}{教授} % 指導教員

% 概要 %%%%%%%%%%%%%%%%%%%%%%%%%%%%%%%%%%%%%%%%%%%%%%%%%%%%%%%%%%%%%%%%%%%%%

\abstract{
ここに内容梗概を書く．

内容梗概の中でも，本文と同様に1行空けで改段落できる．

内容梗概は3段落で書くのがよい．
最初の段落で本研究の目的について述べる．
次の段落で提案手法について述べる．
最後の段落で実装と実験などによって本研究でわかったことについて述べる．

時制は，原則として，最初と次の段落は現在形，最後の段落は過去形になる．

内容梗概は必ず1ページに収めること．
}

\keywords{
キーワードは5個程度とする．「, 」(半角カンマと半角スペース)で区切る．
キーワードはおおむね 社会領域，研究分野 などの広い意味を持つものを先に書き，具体的な技術，手法などの詳細なものを後に書く．
できるかぎり論文タイトルに入っていないものが望ましい．
}

% 本体 %%%%%%%%%%%%%%%%%%%%%%%%%%%%%%%%%%%%%%%%%%%%%%%%%%%%%%%%%%%%%%%%%%%%%%%%%%%%%%%%%%

\begin{document}
\pagenumbering{roman}

\jtitlepage % 表紙
\clearpage % 改ページ
\jabstractpage % 内容梗概
\clearpage % 改ページ
\tableofcontents %目次
\clearpage % 改ページ
\listoftables %表目次
\clearpage % 改ページ
\listoffigures %図目次
\clearpage % 改ページ

\chapter{序論}
\pagenumbering{arabic}

ここには序論を書く．

序論の目的は，自身の行った研究に意義があることを示すことである．
社会的な背景と課題から書き起こし，課題を解決するための工学的なアプローチと
その現状について述べ，自身の提案手法について述べる．最後に本論文の構成を述べる．

序論での図の使用は必要最小限に留める．

\chapter{関連研究(このタイトルは変更すべきかも)}

先行研究を含め，本論文に関係する知識や情報を整理して述べる．

研究に限らず，本論文の内容を理解するうえで必要と考えられる知識の説明を加えてよい．

本章の役割は大きく分けて二つある．
一つ目は，本研究の学術分野内での位置付けを明らかにすることである．
二つ目は，他の研究の紹介を通じて，他の研究では序論で述べた目的を達成できないこと，
よって本研究が必要となることを，具体的に示すことである．
本章では，まず一般的な知識，幅広く関連する研究分野について述べた後，
特に提案するシステムと関連が深い研究について述べる．

なお，各章の冒頭には，その章で何を述べるのかを説明する文章を上記のように記述する．
この部分は5行程度が望ましい．

\section{最初の関連分野}

最初の関連分野について述べる．

関連研究の紹介では，個別の研究をだらだらと羅列するだけにならないように気を付けること．
文章の展開は基本的に大から小である．
関連研究であれば，ある分野に対して，まずその分野の定義や目的を述べ，
簡単に分類し，各分類の特徴と典型的な事例について述べる．
関連研究の詳細まで紹介するのは，本当に自分の研究と関係の深いものだけに留めること．
場合によっては節の末尾で総括を行う．

参考文献\cite{Kameyama92,Favalora02,Tachi15}のテスト．

\chapter{提案手法}

\chapter{実験}

\chapter{考察}

\chapter{結論}

序論に対応させて結論を書く．結論を書く上での注意は以下の通りである．
\begin{itemize}
 \item まず本論文で何を提案し，何を実験し，何を明らかにしたのかを書く．
 \item やり残したことを「今後は，〜する」「今後の課題として〜が挙げられる」と書くと，
       「もう卒業するのにいつやるねん」ということになるので，今後の展望という形で書く．
 \item 序論では社会的な本研究の意義などを書いたはずなので，
       本研究を行ったことによって，どういう社会的な影響を与えることができるか，
       この研究が更に発展したものが社会に普及するとどのようなよいことがあるか，
       といった風呂敷を最後に大きく広げて終わる．
\end{itemize}

% 謝辞 %%%%%%%%%%%%%%%%%%%%%%%%%%%%%%%%%%%%%%%%%%%%%%%%%%%%%%%%%%%%%%%%%%%%%%%%%%%%%%%%%%

\acknowledgements

本研究を進めるにあたり...

% 参考文献 %%%%%%%%%%%%%%%%%%%%%%%%%%%%%%%%%%%%%%%%%%%%%%%%%%%%%%%%%%%%%%%%%%%%%%%%%%%%%%

\bibliography{reference}
\bibliographystyle{junsrt}



% 付録 %%%%%%%%%%%%%%%%%%%%%%%%%%%%%%%%%%%%%%%%%%%%%%%%%%%%%%%%%%%%%%%%%%%%%%%%%%%%%%%%%%

\appendix

\section{本文と直接関係ない知識}

付録は無いことが多いが，どうしても書いておきたいことは書いてもよい．
無い場合は，\verb|\appendix|からまとめて削除する．

% 研究業績

\papers %%%%%%%%%%%%%%%%%%%%%%%%%%%%%%%%%%%%%%%%%%%%%%%%%%%%%%%%%%%%%%%%%%%%%%%%%%%%%%

ここに，あれば自分の業績リスト．無ければ項目ごとコメントアウト．

\end{document}
